\chaves{Data Warehouse; Data Mart; Modelagem Dimensional; ETL; Business Intelligence.} % Palavras chave separadas por ponto e vírgula

% Não deixe linha em branco após \begin{resumo}
\begin{resumo}
Este relatório documenta a construção de um Data Warehouse (DW) Organizacional a partir da integração de dois sistemas transacionais de origem, o Northwind (distribuidora de alimentos) e a Mercearia (comércio varejista), e a derivação de três Data Marts dimensionais a partir dele. O DW Organizacional foi modelado segundo a abordagem de Inmon, orientado a assunto (Tempo, Geografia, Pessoas, Produto e Parceiros, Transações), integrado por meio de identificadores corporativos que resolvem as chaves de cada sistema de origem, não-volátil e variante no tempo, com historização por versionamento (\textit{slowly changing dimensions} do tipo 2) nas entidades cliente, funcionário e produto. A partir desse modelo corporativo foram derivados três Data Marts em esquema estrela, seguindo a abordagem de Kimball e uma arquitetura de barramento com dimensões conformadas: o Mart Vendas, com grão de item vendido, cobrindo os dois sistemas de origem; o Mart Compras, com grão de item comprado, exclusivo da Mercearia; e o Mart Logística e Entregas, com grão de pedido, exclusivo do Northwind, único sistema com o conceito de frete e prazo de entrega. O processo de extração, transformação e carga foi orquestrado com Apache Airflow, com uma carga inicial completa e uma carga incremental baseada em marca d'água sobre o maior identificador de origem já carregado, usando uma área de \textit{staging} intermediária. Um catálogo de metadados documenta a origem, o tipo e a regra de transformação de cada campo do DW, incluindo os dados externos e os campos sem correspondência transacional direta. Por fim, a análise dos dados foi realizada no Metabase, com um painel de indicadores por Data Mart. O ambiente completo foi containerizado com Docker Compose, reunindo os bancos transacionais, o DW em PostgreSQL, o orquestrador de ETL e a ferramenta de \textit{Business Intelligence}.
\end{resumo}
